
\section{Making measurements projective}
\label{sec:making-measurements-projective}

A recurring theme in $\MIP^*$ research is that 
projective measurements are significantly easier to manipulate
than general POVM measurements.
As just one example among many,
when $A = \{A^x_a\}$ and $B = \{B^x_a\}$ are projective measurements,
the ``$\approx_\delta$" and ``$\simeq_\delta$" distances are equivalent.
In other words:
\begin{equation*}
A^x_a \ot I \simeq_{\delta} I \ot B^x_a
\qquad
\iff
\qquad
A^x_a \ot I \approx_{2\delta} I \ot B^x_a.
\end{equation*}
This is Fact~$4.13$ from~\cite{NW19}.
On the other hand when~$A$ and~$B$
are not projective,
one can find examples of measurements
where $A^x_a \ot I \approx_{\delta} I \ot B^x_a$
for $\delta \rightarrow 0$
but one can show $A^x_a \ot I \simeq_{\eps} I \ot B^x_a$
only for $\eps \rightarrow 1$;
this is Remark~$4.15$ from~\cite{NW19}.
As a result, it is important to be able to convert
POVM measurements to projective measurements
whenever possible.

In this section, we will survey two tools for doing so.
The first of these is the textbook Naimark dilation theorem.
In our setting, it states the following.

\begin{theorem}[Naimark dilation]\label{thm:naimark}
Let $\ket{\psi}$ be a state in $\calH_{\mathrm{A}} \ot \calH_{\mathrm{B}}$.
Let $A = \{A^x_a\}$ be a sub-measurement acting on $\calH_{\mathrm{A}}$
and $B=\{B^y_b\}$ be a sub-measurement acting on $\calH_{\mathrm{B}}$.
Then there exists
\begin{enumerate}
\item Hilbert spaces $\calH_{\mathrm{A}_{\mathsf{aux}}}$ and $\calH_{\mathrm{B}_{\mathsf{aux}}}$,
\item a state $\ket{\mathsf{aux}} \in \calH_{\mathrm{A}_{\mathsf{aux}}} \ot \calH_{\mathrm{B}_{\mathsf{aux}}}$,
\item and two measurements  $\widehat{A} = \{\widehat{A}^x_a\}$ and $\widehat{B}=\{\widehat{B}^y_b\}$
acting on $\calH_{\mathrm{A}} \ot \calH_{\mathrm{A}_{\mathsf{aux}}}$
and $\calH_{\mathrm{B}} \ot \calH_{\mathrm{B}_{\mathsf{aux}}}$,
respectively,
\end{enumerate}
such that the following is true.
If we write $\ket{\widehat{\psi}} = \ket{\psi} \otimes \ket{\mathsf{aux}}$,
then for all $x,y, a, b$,
 \begin{equation*}
 \bra{\psi} A^x_a \ot B^y_b \ket{\psi}
 = \bra{\widehat{\psi}} \widehat{A}^x_a \ot \widehat{B}^y_b \ket{\widehat{\psi}}.
 \end{equation*}
 In addition, $\ket{\mathsf{aux}}$ is a \emph{product state},
 meaning that we can write it as
 \begin{equation*}
 \ket{\mathsf{aux}}= \ket{\mathsf{aux}_{\mathrm{A}}} \ot \ket{\mathsf{aux}_{\mathrm{B}}},
 \end{equation*}
for $\ket{\mathsf{aux}_{\mathrm{A}}}$ in $\calH_{\mathrm{A}_{\mathsf{aux}}}$
and $\ket{\mathsf{aux}_{\mathrm{B}}}$ in $\calH_{\mathrm{B}_{\mathsf{aux}}}$.
\end{theorem}
The second of these is the ``orthogonalization lemma" from~\cite{KV11}.
In the setting of symmetric strategies, it states the following.

\begin{theorem}[Orthogonalization lemma]\label{thm:orthonormalization}
Let $\ket{\psi}$ be a permutation-invariant state,
and let $A = \{A_a\}$ be a sub-measurement with strong self-consistency
\begin{equation*}
\sum_a \bra{\psi} A_a \otimes A_a \ket{\psi} \geq \sum_a \bra{\psi} A_a \otimes I \ket{\psi} -\zeta.
\end{equation*}
Then there exists a projective sub-measurement $P = \{P_a\}$ such that
\begin{equation*}
A_a \otimes I \approx_{100 \zeta^{1/4}} P_a \otimes I.
\end{equation*}
\end{theorem}

These two results serve largely the same purpose,
and for our applications it will typically suffice to use either one.
That said, there are tradeoffs when choosing to use one or the other.
Naimark dilation is convenient because it can be applied to any family of sub-measurements
and it preserves measurement outcome probabilities exactly.
On the other hand, it requires swapping out the state $\ket{\psi}$ and measurements~$A$ and~$B$
for $\ket{\widehat{\psi}}$ and $\widehat{A}$ and $\widehat{B}$, respectively,
which can be notationally cumbersome.
Most works opt to skip adding the hats,
and instead say something to the effect of ``using Naimark, we can assume that~$A$ and~$B$ are projective."
However, there is subtlety in doing so,
which is that Naimark dilation does not necessarily preserve ``$\approx_{\delta}$" statements;
in other words, $A^x_a \ot I \approx_{\delta} I \ot B^x_a$
does not necessarily imply $\widehat{A}^x_a \ot I \approx_{\delta} I \ot \widehat{B}^x_a$.
(As an example, \Cref{ex:easy-but-long} below
provides a case in which $A^x_a \ot I \approx_{0} I \ot B^x_a$,
but $\widehat{A}^x_a \ot I \approx_{\delta} I \ot \widehat{B}^x_a$ only holds for $\delta \geq 1$.)
This can lead to trouble if the same notation is used for~$A, B$ and $\widehat{A}, \widehat{B}$,
as ``$\approx_{\delta}$" statements derived before applying Naimark might not necessarily hold after applying Naimark.
(That said, as pointed out at the end of Section~$4.4$ of~\cite{NW19},
``$\approx_{\delta}$" statements are often derived as a consequence of ``$\simeq_{\delta}$" statements.
And as Naimark preserves ``$\simeq_{\delta}$" statements,
one could use them to simply rederive any ``$\approx_{\delta}$" statements post-Naimark.)

The downsides of using the orthogonalization lemma
are that it can only be applied to measurements which are strongly self-consistent
and it introduces additional error.
As we will see below, its proof is significantly more complicated than Naimark dilation.
On the plus side, it has none of the notational baggage that Naimark brings,
and so it is more concrete to use.
In addition, it can lead to stronger results,
as it does not require introducing an auxiliary state $\ket{\mathsf{aux}}$.

In this work, we will opt to use the orthogonalization lemma rather than Naimark dilation.
However, we will include proofs of both for completeness.
We will prove \Cref{thm:naimark} in \Cref{sec:naimark}
and then \Cref{thm:orthonormalization} in \Cref{sec:orthogonalization}.

\subsection{Naimark dilation}\label{sec:naimark}

To prove \Cref{thm:naimark}, we will first need the following lemma.

\begin{lemma}
  \label{lem:naimark-helper}
  Let $A = \{A_{a}\}$ be a sub-measurement with $k$ distinct outcomes $a \in \calA$, and let $\ket{\mathsf{aux}} \in \C^{k+1}$ be any state. Then there exists a projective sub-measurement $\widehat{A} = \{\widehat{A}_{a}\}$ such that
  for each outcome~$a$,
  \[ (I \ot \bra{\mathsf{aux}} ) \cdot \widehat{A}_{a} \cdot (I \ot \ket{\mathsf{aux}}) =
    A_{a}. \]
\end{lemma}
\begin{proof}
We will consider an orthonormal basis for $\C^{k+1}$
consisting of a vector $\ket{a}$ for each $a \in \calA$ and the vector $\ket{\bot}$.
Let~$U$ be any unitary such that for each vector $\ket{\psi} \in \C^d$,
\begin{align*}
U \cdot (\ket{\psi} \ot \ket{\mathsf{aux}})
&= \sum_{a \in \calA} ((A_{a})^{1/2} \ket{\psi}) \ot \ket{a} + ((1 - A)^{1/2}\ket{\psi}) \ot \ket{\bot}\\
&= \Big(\sum_{a \in \calA} (A_{a})^{1/2} \ot \ket{a} + (1 - A)^{1/2} \ot \ket{\bot}\Big)
		\cdot \ket{\psi}.
\end{align*}
This implies that
\begin{equation*}
U \cdot (I \ot \ket{\mathsf{aux}})
= \sum_{a \in \calA} (A_{a})^{1/2} \ot \ket{a} + (1 - A)^{1/2} \ot \ket{\bot}.
\end{equation*}
Hence, for any $a \in \calA$,
\begin{equation*}
(I \ot \bra{a}) \cdot U \cdot (I \ot \ket{\mathsf{aux}})
= (I \ot \bra{a}) \cdot \Big(\sum_{a \in \calA} (A_{a})^{1/2} \ot \ket{a} + (1 - A)^{1/2} \ot \ket{\bot}\Big)
= (A_a)^{1/2}.
\end{equation*}
  Then the desired projective sub-measurement is
  \[ \widehat{A}_a = U^{\dagger} \cdot (I \ot \ket{a} \bra{a}) \cdot U. \]
  This is because
\begin{align*}
 (I \ot \bra{\mathsf{aux}} ) \cdot \widehat{A}_{a} \cdot (I \ot \ket{\mathsf{aux}})
&=(I \ot \bra{\mathsf{aux}} ) \cdot (U^{\dagger} \cdot (I \ot \ket{a} \bra{a}) \cdot U) \cdot (I \ot \ket{\mathsf{aux}})\\
&=(I \ot \bra{\mathsf{aux}} ) \cdot U^{\dagger} \cdot (I \ot \ket{a}) \cdot (I\ot \bra{a}) \cdot U \cdot (I \ot \ket{\mathsf{aux}})\\
&= (A_a)^{1/2} \cdot (A_a)^{1/2} = A_a.
\end{align*}
This completes the proof.
\end{proof}

Now we prove \Cref{thm:naimark}.
\begin{proof}[Proof of \Cref{thm:naimark}]
For each $x$, let $k$ be the number of outcomes in $A^x$,
and let $\ket{\mathsf{aux}_{\mathrm{A},x}}$ be a state of dimensionality $k+1$.
Let $\widetilde{A}^x$ be the sub-measurement guaranteed by \Cref{lem:naimark-helper}.
Define $\ket{\mathsf{aux}_{\mathrm{B},x}}$ and $\widetilde{B}^x$ similarly.
We define
\begin{equation*}
\ket{\widehat{\psi}} = \ket{\psi} \ot \big(\ot_{x} \ket{\mathsf{aux}_{\mathrm{A},x}}\big)
					\ot \big(\ot_{x} \ket{\mathsf{aux}_{\mathrm{B},x}}\big)
\end{equation*}
and
\begin{align*}
\widehat{A}^x_a
&= \widetilde{A}^x_a \ot \big(\ot_{z \neq x} I_{\mathsf{aux}_{\mathrm{A},z}}\big)\\
\widehat{B}^y_b
&= \widetilde{B}^y_b \ot \big(\ot_{z\neq y} I_{\mathsf{aux}_{\mathrm{B},z}}\big).
\end{align*}
Then for all $x,y,z,b$,
\begin{align*}
\bra{\widehat{\psi}} \widehat{A}^x_a \ot \widehat{B}^y_b \ket{\widehat{\psi}}
& = \bra{\psi} \ot \bra{\mathsf{aux}_{\mathrm{A},x}} \ot \bra{\mathsf{aux}_{\mathrm{B},y}}
	\cdot \widetilde{A}^x_a \ot \widetilde{B}^y_b \cdot \ket{\psi} \ot \ket{\mathsf{aux}_{\mathrm{A},x}} \ot \ket{\mathsf{aux}_{\mathrm{B},y}}\\
&= 	\bra{\psi} A^x_a \ot B^y_b \ket{\psi}.
\end{align*}
This completes the proof.
\end{proof}

\begin{example}[Naimark does not preserve ``$\approx_{\delta}$'']\label{ex:easy-but-long}
We now carry out a simple example
in which Naimark preserves ``$\simeq_\delta$'' statements without preserving ``$\approx_\delta$'' statements.
Let $\ket{\psi}$ be an arbitrary state in $\calH_{\mathrm{A}} \ot \calH_{\mathrm{B}}$,
where $\calH_{\mathrm{A}} = \calH_{\mathrm{B}} = \C^d$.
In addition, let $A = \{A_0, A_1\}$ and $B = \{B_0, B_1\}$ be the two-outcome measurements
in which $A_0 = A_1 = B_0 = B_1 = \frac{1}{2} \cdot I_{d \times d}$.
Then
\begin{align*}
\sum_{a \neq b} \bra{\psi} A_a \ot B_b\ket{\psi}
&= \bra{\psi} A_0 \ot B_1\ket{\psi} + \bra{\psi} A_1 \ot B_0\ket{\psi}\\
&= \frac{1}{4} \cdot \bra{\psi} I \ot I \ket{\psi} + \frac{1}{4} \cdot \bra{\psi} I \ot I \ket{\psi}
= \frac{1}{2},
\end{align*}
and so $A_a \ot I \simeq_{1/2} I \ot A_a$.
In addition,
\begin{equation*}
\sum_a \Vert (A_a \ot I - I \ot B_a) \ket{\psi}\Vert^2
= \sum_a \Big\Vert \Big(\frac{1}{2} \cdot I \ot I - \frac{1}{2} \cdot I \ot I\Big) \ket{\psi}\Big\Vert^2
= 0,
\end{equation*}
and so $A_a \ot I \approx_0 I \ot B_a$.

Now we carry out the steps of Naimark dilation.
Because~$A$ and~$B$ are two-outcome,
we will set $\calH_{\mathrm{A}_{\mathsf{aux}}} = \calH_{\mathrm{B}_{\mathsf{aux}}} = \C^2$,
spanned by the basis vectors $\ket{0}$ and $\ket{1}$.
(\Cref{lem:naimark-helper} actually asks that the auxiliary space have dimension $2+1 = 3$,
where the third dimension is spanned by the basis vector $\ket{\bot}$.
However, this is unnecessary for this example because~$A$ and~$B$
are measurements rather than sub-measurements.)
We will choose our two local auxiliary states
to be
\begin{equation*}
\ket{\mathsf{aux}_{\mathrm{A}}} = \ket{\mathsf{aux}_{\mathrm{B}}}
	= \ket{+} = \frac{1}{\sqrt{2}}\ket{0} + \frac{1}{\sqrt{2}}\ket{1}.
\end{equation*}
We first construct $\widehat{A}$ by following the proof of \Cref{lem:naimark-helper}.
To begin, it asks for~$U$ to be a unitary such that for any $\ket{\phi}$ in $\calH_{\mathrm{A}}$,
\begin{align*}
U \cdot (\ket{\phi} \ot \ket{\mathsf{aux}_{\mathrm{A}}})
&=U \cdot (\ket{\phi} \ot \ket{+})\\
&=  ((A_{0})^{1/2} \ket{\phi}) \ot \ket{0} + ((A_{1})^{1/2} \ket{\phi}) \ot \ket{1}\\
&= \Big(\Big(\frac{1}{2}\cdot I\Big)^{1/2} \ket{\phi}\Big) \ot \ket{0}
	+ \Big(\Big(\frac{1}{2}\cdot I\Big)^{1/2} \ket{\phi}\Big) \ot \ket{1}\\
&=\frac{1}{\sqrt{2}} \ket{\phi} \ot \ket{0} + \frac{1}{\sqrt{2}} \ket{\phi} \ot \ket{1}\\
&= \ket{\phi} \ot \ket{+}.
\end{align*}
As a result, we can simply take~$U$ to be the identity matrix.
Thus,
\begin{equation*}
\widehat{A}_a = U^{\dagger} \cdot (I \ot \ket{a} \bra{a}) \cdot U = I \ot \ket{a}\bra{a}.
\end{equation*}
By a similar argument, we can take $\widehat{B}_a = I \ot \ket{a}\bra{a}$ for $a \in \{0, 1\}$ as well.

Now we set $\ket{\widehat{\psi}} = \ket{\psi}\ot \ket{\mathsf{aux}_{\mathrm{A}}} \ot \ket{\mathsf{aux}_{\mathrm{B}}}$.
By \Cref{thm:naimark},
we already know that $\widehat{A}_a \ot I \approx_0 I \ot \widehat{B}_a$
on state $\ket{\widehat{\phi}}$ because this holds for the un-hatted state and measurements.
On the other hand,
we will now show that $\widehat{A}_a \ot I \approx_{\delta} I \ot \widehat{B}_a$
only for $\delta \geq 1$, in contrast to the un-hatted case where it holds for $\delta = 0$.
To see this, we calculate as follows.
\begin{align}
&\sum_a \Vert (\widehat{A}_a \ot I_{\mathrm{B}, \mathsf{aux}_{\mathrm{B}}}
	- I_{\mathrm{A}, \mathsf{aux}_{\mathrm{A}}} \ot \widehat{B}_a) \ket{\widehat{\psi}}\Vert^2\nonumber\\
 ={}& \sum_a \Vert ((I_{\mathrm{A}} \ot \ket{a}\bra{a}) \ot I_{\mathrm{B}, \mathsf{aux}_{\mathrm{B}}}
	- I_{\mathrm{A}, \mathsf{aux}_{\mathrm{A}}} \ot (I_{\mathrm{B}} \ot \ket{a}\bra{a}))
		\ket{\psi} \ot \ket{+} \ot \ket{+}\Vert^2\nonumber\\
={}& \sum_a \Big\Vert \Big(\frac{1}{\sqrt{2}} \ket{\psi} \ot \ket{a}\ot\ket{+}
			- \frac{1}{\sqrt{2}} \ket{\psi} \ot \ket{+} \ot \ket{a}\Big) \Big\Vert^2\nonumber\\
={}& \sum_a \frac{1}{2} \cdot \braket{\psi \mid \psi} 
	\cdot (\bra{a} \ot \bra{+}- \bra{+} \ot \bra{a}) \cdot (\ket{a} \ot \ket{+}- \ket{+} \ot \ket{a}).\label{eq:awefea}
\end{align}
For each~$a$, we have that
\begin{align*}
&(\bra{a} \ot \bra{+}- \bra{+} \ot \bra{a}) \cdot (\ket{a} \ot \ket{+}- \ket{+} \ot \ket{a})\\
={}& 2 - \braket{a \mid +} \cdot \braket{+ \mid a} - \braket{+ \mid a} \cdot \braket{a \mid +}\\
={}& 2 - \frac{1}{\sqrt{2}}\cdot \frac{1}{\sqrt{2}} - \frac{1}{\sqrt{2}} \cdot \frac{1}{\sqrt{2}}\\
={}& 1.
\end{align*}
Plugging this into \Cref{eq:awefea}, we arrive at our bound.
\end{example}

\subsection{Orthogonalization lemma}\label{sec:orthogonalization}

To show \Cref{thm:orthonormalization},
we first show it for the case when~$A$ is a measurement, rather than a sub-measurement.
Our proof of this case will even work in the more general setting
when the strategy is not assumed to be symmetric,
meaning that~$\ket{\psi}$ is not necessarily permutation-invariant and Player~$\mathrm{B}$'s measurement~$B$ may not be equal to Player~$\mathrm{A}$'s measurement.
This is the content of the following lemma, whose proof we defer till later. 

\begin{lemma}[Orthogonalization lemma for measurements]\label{lem:orthonormalization-main-lemma}
Let $\ket{\psi}$ be a state (which is not necessarily permutation-invariant),
and let $A = \{A_a\}$ and $B = \{B_a\}$ be measurements such that
\begin{equation*}
A_a \ot I \simeq_{\zeta} I \ot B_a
\end{equation*}
on state $\ket{\psi}$.
Then there exists a projective sub-measurement $P = \{P_a\}$ such that
\begin{equation*}
A_a \otimes I \approx_{84 \zeta^{1/4}} P_a \otimes I.
\end{equation*}
\end{lemma}

(Note that when $\ket{\psi}$ is permutation-invariant and $B = A$, the condition $A_a \ot I \simeq_{\zeta} I \ot B_a$
is equivalent to~$A_a$ being $\zeta$-strongly self-consistent, by \Cref{prop:other-two-notions-of-self-consistency}.
This is because~$A$ is a measurement.)

\ignore{
\begin{lemma}[Orthogonalization lemma for measurements]\label{lem:orthonormalization-main-lemma}
Let $\ket{\psi}$ be a permutation-invariant state,
and let $A = \{A_a\}$ be a measurement with strong self-consistency
\begin{equation*}
\sum_a \bra{\psi} A_a \otimes A_a \ket{\psi} \geq 1 -\zeta.
\end{equation*}
Then there exists a projective sub-measurement $P = \{P_a\}$ such that
\begin{equation*}
A_a \otimes I \approx_{76 \zeta^{1/4}} P_a \otimes I.
\end{equation*}
\end{lemma}
}

We now prove \Cref{thm:orthonormalization} by reducing the general (sub-measurement) case
to the case when~$A$ is a measurement.

\begin{proof}[Proof of \Cref{thm:orthonormalization} assuming \Cref{lem:orthonormalization-main-lemma}]
Let $A = \{A_a\}$ be a sub-measurement with outcomes $a \in \calA$ whose strong self-consistency is
\begin{equation*}
\sum_a \bra{\psi} A_a \otimes A_a \ket{\psi} \geq \sum_a \bra{\psi} A_a \otimes I \ket{\psi} -\zeta.
\end{equation*}
Note that
\begin{equation*}
\bra{\psi} A \otimes A \ket{\psi}
= \sum_{a, b} \bra{\psi} A_a \otimes A_b \ket{\psi}
\geq \sum_{a} \bra{\psi} A_a \otimes A_a \ket{\psi}
\geq  \bra{\psi} A \otimes I \ket{\psi} -\zeta.
\end{equation*}
Rearranging,
\begin{align*}
\bra{\psi} A \otimes (I - A)\ket{\psi}
&= \bra{\psi} A \otimes I \ket{\psi} - \bra{\psi} A \otimes A \ket{\psi}\\
&\leq (\bra{\psi} A \otimes A \ket{\psi} + \zeta) - \bra{\psi} A \otimes A \ket{\psi}
 = \zeta.
\end{align*}
As a result,
\begin{align*}
\bra{\psi} (I-A) \otimes (I-A) \ket{\psi}
&= \bra{\psi} (I-A) \otimes I \ket{\psi} - \bra{\psi} (I-A) \otimes A \ket{\psi}\\
&\geq \bra{\psi} (I-A) \otimes I \ket{\psi} - \zeta.
\end{align*}


Let $\widehat{A}$ be the POVM measurement with outcomes in $\widehat{\calA} = \calA \cup \{\bot\}$ defined as
\begin{equation*}
\widehat{A}_a
= \left\{\begin{array}{cl}
		A_a & \text{if $a \in \calA$},\\
		(I-A) & \text{if $a = \bot$.}
		\end{array}\right.
\end{equation*}
Then the strong self-consistency of $\widehat{A}$ is
\begin{align*}
\sum_{a \in \widehat{\calA}} \bra{\psi} \widehat{A}_a \otimes \widehat{A}_a \ket{\psi}
& = \sum_{a \in \calA} \bra{\psi} A_a \otimes A_a \ket{\psi}
	+ \bra{\psi} \widehat{A}_{\bot} \otimes \widehat{A}_{\bot} \ket{\psi}\\
& \geq \Big(\sum_{a \in \calA} \bra{\psi} A_a \otimes I \ket{\psi} - \zeta\Big)
	+ \bra{\psi} (I-A) \otimes (I-A) \ket{\psi}\\
& \geq \Big( \bra{\psi} A \otimes I \ket{\psi} - \zeta\Big)
	+ \Big(\bra{\psi} (I-A) \otimes I \ket{\psi} - \zeta\Big)\\
& = 1 - 2\zeta.
\end{align*}
Because $\widehat{A}$ is a measurement,
\Cref{prop:other-two-notions-of-self-consistency} states that this is equivalent to
\begin{equation*}
\widehat{A}_a \ot I \simeq_{2\zeta} I \ot \widehat{A}_a.
\end{equation*}
As a result, \Cref{lem:orthonormalization-main-lemma} implies the existence of a projective sub-measurement
$\widehat{P} = \{\widehat{P}_a\}_{a \in \widehat{\calA}}$ such that 
\begin{equation*}
\widehat{A}_a \ot I \approx_{84\cdot(2\zeta)^{1/4}} \widehat{P}_a \ot I.
\end{equation*}
If we define the projective sub-measurement $P = \{P_a\}_{a \in \calA}$
by $P_a = \widehat{P}_a$ for all $a \in \calA$, then
\begin{equation*}
\sum_{a \in \calA} \Vert (A_a - P_a) \ot I \ket{\psi} \Vert^2
\leq \sum_{a \in \widehat{\calA}} \Vert (\widehat{A}_a - \widehat{P}_a) \ot I \ket{\psi} \Vert^2
\leq 84 \cdot (2\zeta)^{1/4}.
\end{equation*}
Thus, $A_a \ot I \approx_{84 \cdot (2\zeta)^{1/4}} P_a \ot I$.
We conclude the proof by noting that $84 \cdot (2\zeta)^{1/4} \leq 100 \cdot \zeta^{1/4}$
because $84 \cdot (2)^{1/4} \approx 99.89\leq100$.
\end{proof}

Now we prove \Cref{lem:orthonormalization-main-lemma}.

\begin{proof}[Proof of~\Cref{lem:orthonormalization-main-lemma}]
We note that the lemma as stated is trivial when $\zeta > 1/4$.
As a result, we will assume that
\begin{equation}\label{eq:assumption-on-zeta}
\zeta \leq 1/4.
\end{equation}
Let $A = \{A_a\}$ and~$B = \{B_a\}$ be POVM measurements such that
\begin{equation*}
A_a \ot I \simeq_{\zeta} I \ot B_a.
\end{equation*}
\ignore{
By \Cref{prop:two-notions-of-self-consistency},
\begin{equation}\label{eq:A-approx-delta}
A_a \otimes I\approx_{2\zeta} I\otimes A_a.
\end{equation}
}
By \Cref{prop:simeq-for-measurements}, this implies that
\begin{equation*}
\sum_a \bra{\psi}A_a \ot B_a \ket{\psi} \geq 1 - \zeta.
\end{equation*}
Applying Cauchy-Schwarz, we have
\begin{align*}
\sum_a \bra{\psi} A_a \otimes B_a \ket{\psi}
&= \sum_a \bra{\psi} (A_a \ot I) \cdot (I \ot B_a) \ket{\psi}\nonumber\\
&\leq \sqrt{\sum_a \bra{\psi} (A_a)^2 \ot I \ket{\psi}} \cdot \sqrt{\sum_a \bra{\psi} I \ot (B_a)^2 \ket{\psi}}\nonumber\\
& \leq \sqrt{\sum_a \bra{\psi} (A_a)^2 \ot I \ket{\psi}} \cdot 1.
\end{align*}
Taking the square of both sides,
\begin{equation*}
\sum_a \bra{\psi} (A_a)^2 \ot I \ket{\psi} \geq \Big(\sum_a \bra{\psi} A_a \otimes B_a \ket{\psi}\Big)^2 \geq (1-\zeta)^2
\geq 1 - 2\zeta
= \sum_a \bra{\psi} A_a \ot I \ket{\psi} - 2\zeta. \tag{because~$A$ is a measurement}
\end{equation*}
\ignore{
\begin{align}
\sum_a \bra{\psi} A_a \otimes A_a \ket{\psi}
&= \sum_a \bra{\psi} (A_a \ot I) \cdot (I \ot A_a) \ket{\psi}\nonumber\\
&\leq \sqrt{\sum_a \bra{\psi} (A_a)^2 \ot I \ket{\psi}} \cdot \sqrt{\sum_a \bra{\psi} I \ot (A_a)^2 \ket{\psi}}\nonumber\\
& = \sum_a \bra{\psi} (A_a)^2 \ot I \ket{\psi}.\label{eq:As-on-same-side}
\end{align}
\ignore{We claim that
\begin{equation}\label{eq:As-on-same-side}
\sum_a \bra{\psi} A_a \otimes A_a \ket{\psi}
\approx_{\sqrt{2\zeta}} \sum_a \bra{\psi} (A_a)^2 \otimes I \ket{\psi}.
\end{equation}
To show this, we bound the magnitude of the difference.
\begin{align*}
&\Big|\sum_a \bra{\psi} (A_a \otimes I) \cdot (A_a \otimes I - I \otimes A_a) \ket{\psi}\Big|\\
\leq~&
	\sqrt{\sum_a \bra{\psi} (A_a)^2 \otimes I \ket{\psi}}
	\cdot \sqrt{\sum_a \bra{\psi}(A_a \otimes I - I \otimes A_a)^2 \ket{\psi}}\\
\leq~& 1 \cdot \sqrt{2\zeta}. \tag{by \Cref{eq:A-approx-delta}}
\end{align*}}
As a result,
\begin{align*}
\sum_a \bra{\psi} (A_a)^2 \otimes I \ket{\psi}
&\geq \sum_a \bra{\psi} A_a \otimes A_a \ket{\psi}\tag{by \Cref{eq:As-on-same-side}}\\
&\geq \sum_a \bra{\psi} A_a \otimes I \ket{\psi} -\zeta. \tag{by self-consistency of~$A$}
\end{align*}
}
Rearranging,
\begin{equation}\label{eq:A-looks-projective}
\sum_a \bra{\psi} (A_a - (A_a)^2) \otimes I \ket{\psi}
\leq 2\zeta.
\end{equation}
This is all we need the~$B$ measurement for; henceforth, we will derive consequences of~\Cref{eq:A-looks-projective}.

In our next lemma, we convert each~$A_a$ to a projective matrix~$R_a$
by rounding each of $A_a$'s large eigenvalues to~$1$
and small eigenvalues to~$0$.

\begin{lemma}[Rounding to projectors]\label{lem:projective-non-measurement}
There exists a set of projective matrices $\{R_a\}$ such that
\begin{equation*}
A_a \ot I \approx_{2\sqrt{\zeta}} R_a \ot I.
\end{equation*}
and
\begin{equation*}
R:= \sum_a R_a \leq (1+2\sqrt{\zeta}) \cdot I.
\end{equation*}
\end{lemma}
\begin{proof}
For each~$a$, we write the eigendecomposition of $A_a$ as follows: 
\begin{equation*}
A_a = \sum_i \lambda_{a, i} \cdot\ket{u_{a, i}} \bra{u_{a, i}}.
\end{equation*}
Let
\begin{equation}\label{eq:bound-on-delta}
0 < \delta \leq 1/2
\end{equation}
be a number to be decided later.
Let $\mathsf{trunc}_\delta:[0, 1] \rightarrow \{0, 1\}$ be the truncation function defined as
\begin{equation*}
\mathsf{trunc}_\delta(x) =
\left\{\begin{array}{rl}
	1 & \text{if } x \geq 1- \delta,\\
	0 & \text{otherwise.}
	\end{array}\right.
\end{equation*}
Then for each~$a$ we define the matrix $R_a$
\begin{equation*}
R_a = \mathsf{trunc}_\delta(A_a) = \sum_i \mathsf{trunc}_\delta(\lambda_{a, i}) \cdot \ket{u_{a, i}} \bra{u_{a, i}}.
\end{equation*}
To analyze this, we will require the following technical lemma.
\begin{lemma}\label{lem:trunc-inequality}
For any $x \in [0, 1]$,
\begin{equation*}
(x - \mathsf{trunc}_\delta(x))^2 \leq \frac{1}{\delta} \cdot (x - x^2).
\end{equation*}
\end{lemma}
\begin{proof}
This is proved by case analysis.
First, suppose that $x \geq 1 - \delta$.
This implies that $\mathsf{trunc}_\delta(x) = 1$.
In addition, because $\delta \leq 1/2$ by \Cref{eq:bound-on-delta},
$
x \geq 1 - \delta \geq 1/2 \geq \delta.
$
Thus,
\begin{align*}
(x - \mathsf{trunc}_\delta(x))^2
& = (1-x)^2\\
 &\leq (1-x)\\
&\leq (1-x) \cdot \frac{x}{\delta} \tag{because $x \geq \delta$}\\
&= \frac{1}{\delta} \cdot (x - x^2).
\end{align*}
Next, suppose that $x < 1 -\delta$.
This implies that $\mathsf{trunc}_\delta(x) = 0$.
Thus,
\begin{align*}
(x - \mathsf{trunc}_\delta(x))^2
& = x^2\\
 &\leq x\\
&\leq x \cdot \frac{(1-x)}{\delta} \tag{because $x \leq 1- \delta$}\\
&= \frac{1}{\delta} \cdot (x - x^2).
\end{align*}
This concludes the proof.
\end{proof}

As a result, for each~$a$, \Cref{lem:trunc-inequality} implies that
\begin{equation*}
(A_a - R_a)^2 = (A_a - \mathsf{trunc}_\delta(A_a))^2 \leq \frac{1}{\delta} \cdot (A_a - (A_a)^2).
\end{equation*}
Thus,
\begin{align*}
\sum_a \Vert (A_a - R_a) \ot I \ket{\psi}\Vert^2
= \sum_a \bra{\psi} (A_a - R_a)^2 \ot I\ket{\psi}
&\leq \frac{1}{\delta} \cdot \sum_a \bra{\psi} (A_a - (A_a)^2) \ot I \ket{\psi}\\
&\leq \frac{1}{\delta} \cdot 2\zeta. \tag{by \Cref{eq:A-looks-projective}}
\end{align*}
This implies that $A_a \otimes I \approx_{2\zeta/\delta} R_a \otimes I$.

Next, for each $x \in [0, 1]$, it follows from the definition of $\mathsf{trunc}_\delta$ that
\begin{equation*}
\mathsf{trunc}_\delta(x) \leq \left(\frac{1}{1-\delta}\right) \cdot x.
\end{equation*}
Thus, for each~$a$
\begin{equation*}
R_a = \mathsf{trunc}_\delta(A_a) \leq \left(\frac{1}{1-\delta}\right) \cdot A_a.
\end{equation*}
Summing over all~$a$,
\begin{equation*}
R = \sum_a R_a \leq \left(\frac{1}{1-\delta}\right) \cdot \sum_a A_a = \left(\frac{1}{1-\delta}\right) \cdot A
= \left(\frac{1}{1-\delta}\right) \cdot I,
\end{equation*}
because~$A$ is a measurement.
We note that
\begin{equation*}
\frac{1}{1-\delta}
= \frac{1}{1-\delta} \cdot \frac{1 + 2\delta}{1 + 2\delta}
= \frac{1 + 2\delta}{1 + \delta - 2\delta^2}
\leq 1 + 2\delta,
\end{equation*}
because
\begin{equation*}
1 + \delta - 2\delta^2 = 1 + \delta\cdot(1-2\delta)\geq 1
\end{equation*}
when $\delta \leq 1/2$,
which we assumed in \Cref{eq:bound-on-delta}.
Hence,
\begin{equation*}
R \leq (1+ 2\delta) \cdot I.
\end{equation*}


The lemma now follows by setting $\delta = \sqrt{\zeta}$. Note that we required that $\delta$ be at most $1/2$,
which follows from our assumption that $\zeta \leq 1/4$ from \Cref{eq:assumption-on-zeta}.
\ignore{
Finally,
\begin{align*}
&\sum_a R_a (R - R_a) R_a\\
=~&\sum_a R_a R R_a - \sum_a (R_a)^3\\
\leq ~&\left(\frac{1}{1-\delta}\right) \cdot \sum_a R_a A R_a - \sum_a (R_a)^3 \tag{by XXX}\\
\leq ~&\left(\frac{1}{1-\delta}\right) \cdot \sum_a  (R_a)^2 - \sum_a (R_a)^3 \tag{because~$A$ is a sub-measurement}\\
= ~&\left(\frac{1}{1-\delta}\right) \cdot \sum_a  R_a - \sum_a R_a \tag{because the $R_a$'s are projectors}\\
= ~&\left(\frac{1}{1-\delta}\right) \cdot R - R\\
= ~&\left(\frac{1}{1-\delta} -  1\right) \cdot R\\
\leq~& 2\delta \cdot R. \tag{because $\delta \leq 1/2$}
\end{align*}
The lemma now follows by setting $\delta = XXX$. Note that we required that $\delta$ be at most $1/2$,
which follows from our assumption on $\zeta$ from XXX.
}
\end{proof}

Write $d$ for the dimension of the~$A$ matrices.
If the~$\{R_a\}$ matrices from \Cref{lem:projective-non-measurement} formed a projective sub-measurement,
then their total rank would be at most~$d$.
Even if this is not true,
the next lemma shows that we can still post-process them
to reduce their total rank to at most~$d$.

\begin{lemma}[Rank reduction]\label{lem:projective-low-rank-sum}
There exists a set of projection matrices $\{Q_a\}$ such that
\begin{equation*}
A_a \ot I \approx_{12\sqrt{\zeta}} Q_a \ot I.
\end{equation*}
and
\begin{equation*}
Q:= \sum_a Q_a \leq (1+2\sqrt{\zeta}) \cdot I.
\end{equation*}
Furthermore,  $Q$ has bounded total rank:
\begin{equation*}
\sum_a \mathrm{rank}(Q_a) \leq d.
\end{equation*}
\end{lemma}
\begin{proof}
Let $\{R_a\}$ be the set of projective matrices given by \Cref{lem:projective-non-measurement}.
For each $a$, let $r_a$ be the rank of $R_a$.
Let $r = \sum_a r_a$.
If $r \leq d$, then the lemma is trivially satisfied by taking $Q$ to be~$R$ and applying \Cref{lem:projective-non-measurement}.
Henceforth, we will assume that $r > d$.
We want to reduce~$r$ so that it is at most~$d$.
Fortunately, it is already not too much larger than~$d$:
\begin{equation}\label{eq:bound-on-r}
r
= \sum_{a} r_a
= \sum_{a} \trace(R_a)
= \trace(R)
\leq (1 + 2\sqrt{\zeta}) \cdot \trace(I)
= (1+2\sqrt{\zeta}) \cdot d.
\end{equation}
Let $\ket{v_{a, 1}}, \ldots, \ket{v_{a, r_a}}$ be an orthonormal basis for the range of $R_a$, so that
\begin{equation*}
R_a = \sum_{i=1}^{r_a} \ket{v_{a, i}}\bra{v_{a,i}}.
\end{equation*}
To reduce~$r$,
we will throw out those $\ket{v_{a,i}}$'s whose overlap with $\ket{\psi}$ is small.
For each $a$ and $1 \leq i \leq r_a$, we denote the overlap of $\ket{v_{a,i}}$ and $\ket{\psi}$ by
\begin{equation*}
o_{a, i} = \bra{\psi} \cdot (\ket{v_{a,i}}\bra {v_{a,i}} \ot I) \cdot\ket{\psi}.
\end{equation*}
The total overlap is given by
\begin{align}
\sum_a \sum_{i=1}^{r_a} o_{a, i}
&= \sum_a \sum_{i=1}^{r_a} \bra{\psi} \cdot (\ket{v_{a,i}}\bra {v_{a,i}} \ot I) \cdot\ket{\psi}\nonumber\\
&= \sum_a \bra{\psi} R_a \ot I \ket{\psi}\nonumber\\
&= \bra{\psi} R \ot I \ket{\psi}\nonumber\\
&\leq (1+2\sqrt{\zeta}) \cdot \bra{\psi} I \ot I \ket{\psi}\nonumber\\
&= 1+2\sqrt{\zeta}.\label{eq:total-overlap}
\end{align}
Now we define $\mathsf{Large}$ to be the set of large overlaps:
\begin{equation*}
\mathsf{Large} = \{(a, i) \mid \text{$o_{a, i}$ is among the~$d$ largest of the $o_{b, j}$'s}\},
\end{equation*}
where we break ties arbitrarily,
and we define $\mathsf{Small}$ to be the set containing the remaining $(a, i)$'s.
We note that $\mathsf{Large}$ is well-defined and has size~$d$ because $r > d$.
Hence, \Cref{eq:bound-on-r} implies that
\begin{equation}\label{eq:size-of-small-is-small}
|\mathsf{Small}|
= r - |\mathsf{Large}| 
= r - d
\leq (1 + 2\sqrt{\zeta}) \cdot d - d
= 2\sqrt{\zeta} \cdot d
\leq 2\sqrt{\zeta} \cdot r.
\end{equation}
Thus, the small $(a,i)$'s have small total overlap:
\begin{align}
\sum_{(a, i) \in \mathsf{Small}} o_{a, i}
\leq \frac{|\mathsf{Small}|}{r} \sum_{a, i} o_{a,i}
&\leq 2\sqrt{\zeta} \cdot \sum_{a, i} o_{a,i} \tag{by \Cref{eq:size-of-small-is-small}}\\
&\leq 2\sqrt{\zeta} \cdot (1+2\sqrt{\zeta}) \tag{by \Cref{eq:total-overlap}}\\
& \leq 4\sqrt{\zeta}, \label{eq:small-overlaps}
\end{align}
where the final step uses the assumption that $\zeta \leq 1/4$ from \Cref{eq:assumption-on-zeta}.

For each~$a$ we let $\mathsf{Large}_a$ to be set of $i$'s such that $(a,i)$ is contained in $\mathsf{Large}$,
and we define $\mathsf{Small}_a$ similarly.
We define the matrix
\begin{equation*}
Q_a = \sum_{i\in \mathsf{Large}_a} \ket{v_{a, i}}\bra{v_{a,i}}.
\end{equation*}
Then clearly
\begin{equation*}
\sum_a \mathrm{rank}(Q_a)
= \sum_a |\mathsf{Large}_a|
= |\mathsf{Large}|
\leq d.
\end{equation*}
We can compute the difference
\begin{equation*}
R_a - Q_a
= \sum_{i=1}^{r_a} \ket{v_{a, i}}\bra{v_{a,i}} - \sum_{i \in \mathsf{Large}_a}  \ket{v_{a, i}}\bra{v_{a,i}}
= \sum_{i\in \mathsf{Small}_a} \ket{v_{a, i}}\bra{v_{a,i}}.
\end{equation*}
This is a projective Hermitian matrix, which implies that $Q_a \leq R_a$.
As a result,
\begin{equation*}
Q = \sum_a Q_a \leq \sum_a R_a = R \leq (1+2\sqrt{\zeta}) \cdot I.
\end{equation*}
In addition,
\begin{align*}
\sum_a \Vert (R_a - Q_a) \ot I \ket{\psi} \Vert^2
&= \sum_a \bra{\psi}(R_a - Q_a)^2 \ot I \ket{\psi}\\
&= \sum_a \bra{\psi}(R_a - Q_a)\ot I \ket{\psi} \tag{because $R_a - Q_a$ is a projector}\\
&= \sum_a \sum_{i \in \mathsf{Small}_a} \bra{\psi} \cdot (\ket{v_{a,i}}\bra {v_{a,i}} \ot I) \cdot\ket{\psi}\\
&= \sum_a \sum_{i \in \mathsf{Small}_a} o_{a, i}\\
&\leq 4\sqrt{\zeta}. \tag{by \Cref{eq:small-overlaps}}
\end{align*}
This means that
\begin{equation*}
R_a \ot I \approx_{4\sqrt{\zeta}} Q_a \ot I.
\end{equation*}
Since we know that $A_a \ot I \approx_{2\sqrt{\zeta}} R_a \ot I$ by \Cref{lem:projective-non-measurement},
\Cref{prop:triangle-inequality-for-approx_delta} implies that
\begin{equation*}
A_a \ot I \approx_{12\sqrt{\zeta}} Q_a \ot I.
\end{equation*}
\ignore{
Finally,
\begin{align*}
&\sum_a Q_a (Q - Q_a) Q_a\\
=~&\sum_a Q_a Q Q_a - \sum_a (Q_a)^3\\
\leq ~&XXX \cdot \sum_a Q_a A Q_a - \sum_a (Q_a)^3 \tag{by XXX}\\
\leq ~&XXX \cdot \sum_a  (Q_a)^2 - \sum_a (Q_a)^3 \tag{because~$A$ is a sub-measurement}\\
= ~&XXX \cdot \sum_a  Q_a - \sum_a Q_a \tag{because the $Q_a$'s are projectors}\\
= ~&XXX \cdot Q - Q\\
= ~&\left(XXX -  1\right) \cdot Q\\
\leq~& 2\delta \cdot Q.
\end{align*}
}
This completes the proof.
\end{proof}

Henceforth, we let $Q = \{Q_a\}$ be the set of projective matrices given by \Cref{lem:projective-low-rank-sum}.
We now derive a few properties of~$Q$ that follow as a consequence of \Cref{lem:projective-low-rank-sum}.
To begin, we show that~$Q$ is almost as complete as~$A$.

\begin{lemma}[Completeness of~$Q$]\label{lem:Q-completeness}
\begin{equation*}
\bra{\psi} Q \otimes I \ket{\psi}
\geq 1 - 11 \zeta^{1/4}.
\end{equation*}
\end{lemma}
\begin{proof}
To begin, we claim that
\begin{align}
\bra{\psi} Q \otimes I \ket{\psi}
& = \sum_a \bra{\psi} Q_a \otimes I \ket{\psi}\nonumber\\
& = \sum_a \bra{\psi} (Q_a)^2 \otimes I \ket{\psi}\tag{because the $Q_a$'s are projective}\\
& \approx_{5 \zeta^{1/4}} \sum_a \bra{\psi} (Q_a \cdot A_a) \otimes I \ket{\psi}. \label{eq:Q-for-an-A}
\end{align}
To show this, we bound the magnitude of the difference.
\begin{align*}
&\Big|\sum_a \bra{\psi} (Q_a \otimes I) \cdot ((Q_a - A_a) \otimes I) \ket{\psi}\Big|\\
\leq~&\sqrt{\sum_a\bra{\psi} (Q_a)^2 \otimes I \ket{\psi}}
	\cdot \sqrt{\sum_a \bra{\psi} (Q_a - A_a)^2 \otimes I \ket{\psi}}\\
\leq~&  \sqrt{1+2\sqrt{\zeta}} \cdot \sqrt{12\sqrt{\zeta}} \tag{by \Cref{lem:projective-low-rank-sum}}\\
\leq~& \sqrt{2} \cdot \sqrt{12\sqrt{\zeta}},
\end{align*}
where the last line uses the assumption that $\zeta \leq 1/4$ from \Cref{eq:assumption-on-zeta}.
Next, we claim that
\begin{equation}\label{eq:another-Q-for-A}
\eqref{eq:Q-for-an-A}
= \sum_a \bra{\psi} (Q_a \cdot A_a) \otimes I \ket{\psi}
\approx_{4 \zeta^{1/4}} \sum_a \bra{\psi} (A_a)^2 \otimes I \ket{\psi}.
\end{equation}
To show this, we bound the magnitude of the difference.
\begin{align*}
&\Big|\sum_a \bra{\psi} ((Q_a - A_a) \otimes I) \cdot (A_a \otimes I) \ket{\psi}\Big|\\
\leq~&\sqrt{\sum_a \bra{\psi} (Q_a - A_a)^2 \otimes I \ket{\psi}}
	\cdot \sqrt{\sum_a\bra{\psi} (A_a)^2 \otimes I \ket{\psi}}\\
\leq~& \sqrt{12 \sqrt{\zeta}} \cdot 1. \tag{by \Cref{lem:projective-low-rank-sum}}
\end{align*}
In conclusion,
\begin{align*}
\bra{\psi} Q \otimes I \ket{\psi}
& \geq \sum_a \bra{\psi} (A_a)^2 \otimes I \ket{\psi} - 9 \zeta^{1/4} \tag{by \Cref{eq:Q-for-an-A,eq:another-Q-for-A}}\\
& \geq \sum_a \bra{\psi} A_a \otimes I \ket{\psi} - 2\zeta - 9\zeta^{1/4} \tag{by \Cref{eq:A-looks-projective}}\\
& \geq \bra{\psi} A \otimes I \ket{\psi} - 11 \zeta^{1/4} \nonumber\\
& = 1 - 11 \zeta^{1/4}.
\end{align*}
This completes the proof.
\end{proof}

We will also need the following bound on the completeness of the \emph{square root} of~$Q$.
Note that such a bound follows trivially from \Cref{lem:Q-completeness}
when~$Q$ is a sub-measurement because $\sqrt{Q} \geq Q$ when $Q \leq I$.

\begin{lemma}[Completeness of~$\sqrt{Q}$]\label{lem:sqrt-Q-completeness}
\begin{equation*}
\bra{\psi} \sqrt{Q} \otimes I \ket{\psi}
\geq 1 - 12 \zeta^{1/4}.
\end{equation*}
\end{lemma}
\begin{proof}
Let $Q = \sum_i \nu_i \ket{u_i} \bra{u_i}$ be the eigendecomposition of~$Q$.
Then because $Q \leq (1 + 2\sqrt{\zeta}) \cdot I$, each eigenvalue $\nu_i$ is at most $1+2\sqrt{\zeta}$.
Thus,
\begin{equation*}
\sqrt{Q}
= \sum_i \sqrt{\nu_i} \ket{u_i}\bra{u_i}
\geq \frac{1}{\sqrt{1 + 2\sqrt{\zeta}}} \cdot \sum_i \nu_i \ket{u_i}\bra{u_i}
= \frac{1}{\sqrt{1 + 2\sqrt{\zeta}}} \cdot Q.
\end{equation*}
We note that
\begin{equation*}
\frac{1}{\sqrt{1 + 2\sqrt{\zeta}}}
\geq \frac{1}{\sqrt{1 + 2\sqrt{\zeta} + \zeta}}
= \frac{1}{1 + \sqrt{\zeta}}
= \frac{1}{1 + \sqrt{\zeta}} \cdot \Big(\frac{1-\sqrt{\zeta}}{1-\sqrt{\zeta}}\Big)
= \frac{1-\sqrt{\zeta}}{1 - \zeta}
\geq 1 - \sqrt{\zeta}.
\end{equation*}
Hence, $\sqrt{Q} \geq (1-\sqrt{\zeta}) \cdot Q$.
As a result, \Cref{lem:Q-completeness} implies that
\begin{equation*}
\bra{\psi} \sqrt{Q} \otimes I \ket{\psi}
\geq (1-\sqrt{\zeta}) \cdot\bra{\psi} Q \otimes I \ket{\psi}
\geq (1-\sqrt{\zeta}) \cdot \left(1 - 11 \zeta^{1/4}\right)
\geq 1 -\sqrt{\zeta} - 11 \zeta^{1/4}.
\end{equation*}
This concludes the proof.
\end{proof}

Finally, we show the following lemma,
which quantifies a sense in which~$Q$ is ``almost projective".

\begin{lemma}[$Q$ is almost projective]\label{lem:q-almost-projective}
\begin{equation*}
\sum_a (Q_a \cdot Q  \cdot Q_a - Q_a) \leq  4\sqrt{\zeta} \cdot I.
\end{equation*}
\end{lemma}
\begin{proof}
\Cref{lem:projective-low-rank-sum} implies that $Q \leq (1 + 2\sqrt{\zeta}) \cdot I$. As a result,
\begin{align*}
&\sum_a Q_a \cdot Q  \cdot Q_a - \sum_a Q_a\\
\leq ~&(1 + 2\sqrt{\zeta}) \cdot \sum_a Q_a \cdot I  \cdot Q_a - \sum_a Q_a \\
= ~&(1+2\sqrt{\zeta}) \cdot \sum_a  Q_a - \sum_a Q_a \tag{because the $Q_a$'s are projectors}\\
= ~&(1+2\sqrt{\zeta}) \cdot Q - Q\\
= ~& 2\sqrt{\zeta} \cdot Q\\
\leq~& 2\sqrt{\zeta} \cdot (1 + 2\sqrt{\zeta}) \cdot I\\
\leq~& 2 \sqrt{\zeta} \cdot 2 \cdot I. \tag{by \Cref{eq:assumption-on-zeta}}
\end{align*}
This completes the proof.
\end{proof}

We now arrive at the most important definition in this proof,
which is a natural matrix decomposition for the $Q_a$ matrices.

\begin{definition}[Matrix decomposition of~$Q_a$]
For each $a$, let $m_a$ be the rank of $Q_a$.
Let $\ket{v_{a, 1}}, \ldots, \ket{v_{a, m_a}}$ be an orthonormal basis for the range of $Q_a$, so that
\begin{equation*}
Q_a = \sum_{i=1}^{m_a} \ket{v_{a, i}}\bra{v_{a,i}}.
\end{equation*}
Let $m = \sum_a m_a$, and consider an orthonormal basis of $\C^m$ consisting of vectors $\ket{a, i}$ for each $a$ and $1 \leq i \leq m_a$.
For each $a$, define the matrix
\begin{equation*}
X_a = \sum_{i=1}^{m_a} \ket{a, i} \bra{v_{a, i}}.
\end{equation*}
In addition, define the matrix
\begin{equation}\label{eq:looks-like-singular-value-decomposition}
X = \sum_a X_a =  \sum_a \sum_{i=1}^{m_a} \ket{a, i} \bra{v_{a, i}}.
\end{equation}
Finally, we let $T = \{T_a\}$ be the projective measurement on $\C^m$ defined as
\begin{equation*}
T_a = \sum_{i=1}^{m_a} \ket{a, i}\bra{a, i}.
\end{equation*}
\end{definition}

The next pair of lemmas will prove some basic properties of the~$X_a$ matrices.

\begin{lemma}\label{lem:xa-t}
For each~$a$, $\displaystyle
X_a = T_a \cdot X.
$
\end{lemma}
\begin{proof}
This is a simple calculation:
\begin{equation*}
T_a \cdot X = \Big(\sum_i^{m_a} \ket{a, i} \bra{a,i}\Big)\cdot \Big( \sum_a \sum_{i=1}^{m_a}  \ket{a, i}\bra{v_{a, i}}\Big) 
= \sum_{i=1}^{m_a} \ket{a, i}\bra{v_{a, i}} = X_a.\qedhere
\end{equation*}
\end{proof}

\begin{lemma}[$Q_a$ restated]\label{lem:qa-restated}
For each~$a$,
\begin{equation*}
Q_a = X^\dagger_a \cdot X_a = X^\dagger \cdot T_a \cdot X = X_a^\dagger \cdot X.
\end{equation*}
\end{lemma}
\begin{proof}
The first equality follows from
\begin{align*}
X_a^\dagger \cdot X_a
&= \Big(\sum_{i=1}^{m_a} \ket{v_{a, i}}\bra{a, i} \Big) \cdot \Big(\sum_{j=1}^{m_a} \ket{v_{a, j}}\bra{a, j} \Big)^\dagger\nonumber\\
&= \sum_{i, j=1}^{m_a} \ket{v_{a, i}}\bra{a, i} \cdot \ket{a, j} \bra{v_{a, j}}
= \sum_{i=1}^{m_a} \ket{v_{a,i}}\bra{v_{a,i}} = Q_a.
\end{align*}
The remaining equalities follow from \Cref{lem:xa-t} and the fact that~$T$ is a projective measurement.
\begin{equation*}
X_a^\dagger \cdot X_a
= (X^\dagger \cdot T_a) \cdot (T_a  \cdot X)
= X^\dagger \cdot T_a \cdot X = X_a^\dagger \cdot X.\qedhere
\end{equation*}
\end{proof}

Now we introduce our main tool for studying~$X$, which is via its singular value decomposition.

\begin{definition}[SVD of~$X$]
Let $X = U \cdot \Sigma_{m \times d} \cdot V^\dagger$ be the singular value decomposition (SVD) of~$X$.
Because~$X$ is an $m \times d$ matrix,
the definition of the SVD states that
$U$ is an $m \times m$ unitary matrix,
$V$ is a $d \times d$ unitary matrix,
and $\Sigma_{m\times d}$ is an $m \times d$ diagonal matrix
with nonnegative real numbers on its diagonal.
\end{definition}

\begin{notation}
For positive integers $h$ and~$w$,
we will write $I_{h \times w}$ for the $h \times w$ matrix with $1$'s on its diagonal and $0$'s everywhere else.

For integers $h, w \geq m$, we also write $\Sigma_{h \times w}$
for the $h \times w$ diagonal matrix whose diagonal agrees with $\Sigma_{m \times d}$'s;
namely, $(\Sigma_{h \times w})_{i, i} = (\Sigma_{m \times d})_{i,i}$ for all $1 \leq i \leq m$
and $(\Sigma_{h \times w})_{i, i} = 0$ for all $i > m$.
We note that because $\Sigma_{m \times d}$ is a real-valued diagonal matrix,
$(\Sigma_{h \times w})^\dagger = \Sigma_{w \times h}$.
\end{notation}

With this definition,
we can give a helpful expression for the square of~$X$.

\begin{lemma}[$X$ squared]\label{lem:X-squared}
\begin{equation*}
X \cdot X^\dagger = U  \cdot (\Sigma_{m \times m})^2 \cdot U^\dagger,
\quad
\text{and}
\quad
X^\dagger \cdot X = Q = V \cdot (\Sigma_{d \times d})^2 \cdot V^\dagger.
\end{equation*}
\end{lemma}
\begin{proof}
First,
\begin{align*}
X \cdot X^\dagger
&= (U  \cdot \Sigma_{m \times d} \cdot V^\dagger) \cdot (U  \cdot \Sigma_{m \times d} \cdot V^\dagger)^\dagger\\
&= U  \cdot \Sigma_{m \times d} \cdot V^\dagger \cdot V \cdot \Sigma_{d \times m} \cdot U^\dagger\\
&= U  \cdot \Sigma_{m \times d} \cdot I_{d \times d} \cdot \Sigma_{d \times m} \cdot U^\dagger \tag{because~$V$ is a $d \times d$ unitary}\\
&= U  \cdot (\Sigma_{m \times m})^2 \cdot U^\dagger \tag{because~$m \leq d$}.
\end{align*}
Second,
\begin{align*}
X^\dagger \cdot X
& = \sum_a X^\dagger \cdot T_a X  \tag{because~$T$ is a measurement}\\
&= \sum_a Q_a \tag{by \Cref{lem:qa-restated}}\\
&= Q.
\end{align*}
In addition, we can rewrite $X^\dagger \cdot X$ as
\begin{align*}
X^\dagger \cdot X
&= (U  \cdot \Sigma_{m \times d} \cdot V^\dagger)^\dagger \cdot (U  \cdot \Sigma_{m \times d} \cdot V^\dagger)\\
&= V \cdot \Sigma_{d \times m} \cdot U^\dagger \cdot U \cdot \Sigma_{m \times d} \cdot V^\dagger\\
&= V \cdot \Sigma_{d \times m} \cdot I_{m \times m} \cdot \Sigma_{m \times d} \cdot V^\dagger \tag{because~$U$ is an $m \times m$ unitary}\\
&= V \cdot (\Sigma_{d \times d})^2 \cdot V^\dagger.
\end{align*}
This completes the proof.
\end{proof}

The following lemma relates an expression in the $X_a$'s
with an expression in the~$Q_a$'s that appeared previously in \Cref{lem:q-almost-projective}.

\begin{lemma}\label{lem:X-expression-to-Q-expression}
For each~$a$,
\begin{equation*}
X_a^\dagger \cdot (X \cdot X^\dagger - I_{m \times m})^2 \cdot X_a
= Q_a \cdot Q \cdot Q_a - Q_a.
\end{equation*}
\end{lemma}
\begin{proof}
Using~\Cref{lem:X-squared} and~\Cref{lem:qa-restated},
\begin{equation*}
X_a^\dagger \cdot (X \cdot X^\dagger \cdot X \cdot X^\dagger) \cdot X_a
= X_a^\dagger \cdot X \cdot Q \cdot X^\dagger \cdot X_a
= Q_a \cdot Q \cdot Q_a.
\end{equation*}
Similarly, because $Q_a$ is projective,
\begin{equation*}
X_a^\dagger \cdot (X \cdot X^\dagger) \cdot X_a
= (X_a^\dagger \cdot X) \cdot (X^\dagger \cdot X_a)
= Q_a \cdot Q_a
= Q_a.
\end{equation*}
Finally,
\begin{equation*}
X_a^\dagger \cdot (I_{m \times m}) \cdot X_a
= X_a^\dagger \cdot X_a = Q_a.
\end{equation*}
Putting these together,
\begin{align*}
&X_a^\dagger \cdot (X \cdot X^\dagger - I_{m \times m})^2 \cdot X_a\\
 =~& X_a^\dagger \cdot (X \cdot X^\dagger \cdot X \cdot X^\dagger  - 2 \cdot X\cdot X^\dagger + I_{m \times m}) \cdot X_a\\
    =~& Q_a \cdot Q \cdot Q_a  - 2 \cdot Q_a  + Q_a\\
    =~& Q_a \cdot Q \cdot Q_a  -  Q_a.
\end{align*}
This completes the proof.
\end{proof}

Now we are ready to state the projective sub-measurement~$P$
which should approximate~$A$.
Before doing so, we give some intuition for the construction.

\begin{remark}
Suppose that the $Q_a$'s actually formed a projective measurement.
This would imply that the vectors $\ket{v_{a, i}}$, over all $a$ and $1 \leq i \leq m_a$ form an orthonormal set.
Then the SVD would actually have already been provided in \Cref{eq:looks-like-singular-value-decomposition};
for each $a$ and $1 \leq i \leq m_a$,
the corresponding singular value would be~$1$ and the corresponding left- and right-singular vectors would be $\ket{a, i}$ and $\ket{v_{a, i}}$, respectively.
In particular, we would have $\Sigma = I_{m \times d}$ and $X = U \cdot I_{m \times d}\cdot V^\dagger$.
\end{remark}

In reality, we don't know that $\Sigma = I_{m \times d}$.
However, we will construct~$P = \{P_a\}$ as if it were.
This motivates the following definition.

\begin{definition}[Definition of~$P$]
Define the matrix 
\begin{equation*}
\widehat{X} = U  \cdot I_{m \times d} \cdot V^\dagger.
\end{equation*}
In addition, for each $a$, define the matrices
\begin{equation*}
\widehat{X}_a  = T_a \cdot \widehat{X},
\quad P_a  = \widehat{X}_a^\dagger \cdot \widehat{X}_a.
\end{equation*}
\end{definition}

We now give analogues of \Cref{lem:qa-restated,lem:X-squared} for the~$P$ matrices.

\begin{lemma}[$P_a$ restated]\label{lem:pa-restated}
For each~$a$,
\begin{equation*}
P_a = \widehat{X}^\dagger \cdot T_a \cdot \widehat{X} = \widehat{X}_a^\dagger \cdot \widehat{X}.
\end{equation*}
\end{lemma}
\begin{proof}
This follows from the definition of $\widehat{X}^\dagger_a$ and the fact that $T$ is a projective measurement:
\begin{equation*}
P_a
= \widehat{X}_a^\dagger \cdot \widehat{X}_a
= (\widehat{X}^\dagger \cdot T_a) \cdot (T_a  \cdot \widehat{X})
= \widehat{X}^\dagger \cdot T_a \cdot \widehat{X} = \widehat{X}_a^\dagger \cdot \widehat{X}.\qedhere
\end{equation*}
\end{proof}

\begin{lemma}[$\widehat{X}$ squared]\label{lem:X-hat-squared}
\begin{equation*}
\widehat{X} \cdot \widehat{X}^\dagger = I_{m \times m}.
\ignore{
\quad
\text{and}
\quad
\widehat{X}^\dagger \cdot \widehat{X} = V \cdot I_{d \times m} \cdot I_{m \times d} \cdot V^\dagger.
}
\end{equation*}
\end{lemma}
\begin{proof}
First,
\begin{align*}
\widehat{X} \cdot \widehat{X}^\dagger
&= (U  \cdot I_{m \times d} \cdot V^\dagger) \cdot (U  \cdot I_{m \times d} \cdot V^\dagger)^\dagger\\
&= U  \cdot I_{m \times d} \cdot V^\dagger \cdot V \cdot I_{d \times m} \cdot U^\dagger\\
&= U  \cdot I_{m \times d} \cdot I_{d \times d} \cdot I_{d \times m} \cdot U^\dagger \tag{because~$V$ is a $d \times d$ unitary}\\
&= U  \cdot I_{m \times m} \cdot U^\dagger \tag{because~$m \leq d$}\\
&= I_{m \times m},
\end{align*}
where the last step uses the fact that~$U$ is an $m \times m$ unitary.
\ignore{
Second,
\begin{align*}
\widehat{X}^\dagger \cdot \widehat{X}
& = (U  \cdot I_{m \times d} \cdot V^\dagger)^\dagger \cdot (U  \cdot I_{m \times d} \cdot V^\dagger)\\
& = V \cdot I_{d \times m} \cdot U^\dagger \cdot U  \cdot I_{m \times d} \cdot V^\dagger\\
& = V \cdot I_{d \times m} \cdot I_{m \times m}  \cdot I_{m \times d} \cdot V^\dagger \tag{because~$U$ is an $m \times m$ unitary}\\
& = V \cdot I_{d \times m} \cdot I_{m \times d} \cdot V^\dagger.
\end{align*}
We note that $m \leq d$, and so $I_{d \times m} \cdot I_{m \times d}$ is not necessarily equal to $I_{d \times d}$.
}
\end{proof}

Finally, we show two lemmas on quantities involving both $X$ and~$\widehat{X}$.

\begin{lemma}[$X$ times $\widehat{X}$]\label{lem:X-times-X-hat}
\begin{equation*}
X \cdot \widehat{X}^\dagger = U  \cdot \Sigma_{m \times m} \cdot U^\dagger,
\quad
\text{and}
\quad
X^\dagger \cdot \widehat{X} = \sqrt{Q}.
\end{equation*}
\end{lemma}
\begin{proof}
First,
\begin{align*}
X \cdot \widehat{X}^\dagger
&= (U  \cdot \Sigma_{m \times d} \cdot V^\dagger) \cdot (U  \cdot I_{m \times d} \cdot V^\dagger)^\dagger\\
&= U  \cdot \Sigma_{m \times d} \cdot V^\dagger \cdot V \cdot I_{d \times m} \cdot U^\dagger\\
&= U  \cdot \Sigma_{m \times d} \cdot I_{d \times d} \cdot I_{d \times m} \cdot U^\dagger \tag{because~$V$ is a $d \times d$ unitary}\\
&= U  \cdot \Sigma_{m \times m} \cdot U^\dagger \tag{because~$m \leq d$}.
\end{align*}
Second,
\begin{align*}
X^\dagger \cdot \widehat{X}
&= (U  \cdot \Sigma_{m \times d} \cdot V^\dagger)^\dagger \cdot (U  \cdot I_{m \times d} \cdot V^\dagger)\\
&= V \cdot \Sigma_{d \times m} \cdot U^\dagger \cdot U \cdot I_{m \times d} \cdot V^\dagger\\
&= V \cdot \Sigma_{d \times m} \cdot I_{m \times m} \cdot I_{m \times d} \cdot V^\dagger \tag{because~$U$ is an $m \times m$ unitary}\\
&= V \cdot \Sigma_{d \times d} \cdot V^\dagger\\
&= \sqrt{V \cdot (\Sigma_{d \times d})^2 \cdot V^\dagger}\\
&= \sqrt{Q}. \tag{by \Cref{lem:X-squared}}
\end{align*}
This completes the proof.
\end{proof}

\begin{lemma}[Squared difference]\label{lem:squared-difference}
\begin{equation*}
(X - \widehat{X}) \cdot (X - \widehat{X})^\dagger \leq (X \cdot X^\dagger - I_{m \times m})^2.
\end{equation*}
\end{lemma}
\begin{proof}
By~\Cref{lem:X-squared,lem:X-hat-squared,lem:X-times-X-hat},
\begin{align*}
(X - \widehat{X}) \cdot (X - \widehat{X})^\dagger
& = X \cdot X^\dagger  - X\cdot \widehat{X}^\dagger - \widehat{X} \cdot X^\dagger + \widehat{X} \cdot \widehat{X}^\dagger\\
& = U \cdot \Sigma_{m \times m}^2 \cdot U^\dagger  - 2 \cdot U \cdot \Sigma_{m \times m} \cdot U^\dagger+ I_{m \times m}\\
& = U \cdot \left(\Sigma_{m \times m}^2   - 2 \cdot  \Sigma_{m \times m} +  I_{m \times m} \right) \cdot U^\dagger\\
& = U \cdot \left(\Sigma_{m \times m}  -  I_{m \times m}\right)^2 \cdot U^\dagger.
\end{align*}
Because $\Sigma_{m \times m}$ and $I_{m \times m}$ are commuting,
\begin{equation*}
(\Sigma_{m \times m} - I_{m \times m})^2
\leq (\Sigma_{m \times m} + I_{m \times m})^2 \cdot (\Sigma_{m \times m} - I_{m \times m})^2
= (\Sigma_{m\times m}^2 - I_{m \times m})^2.
\end{equation*}
As a result,
\begin{align*}
U \cdot \left(\Sigma_{m \times m}  -  I_{m \times m}\right)^2 \cdot U^\dagger
& \leq U \cdot \left(\Sigma_{m \times m}^2 -  I_{m \times m}\right)^2 \cdot U^\dagger\\
& = \left(U \cdot \Sigma_{m \times m}^2 \cdot U^\dagger - I_{m \times m} \right)^2\\
&= (X \cdot X^\dagger - I_{m \times m})^2. \tag{by \Cref{lem:X-squared}}
\end{align*}
This completes the proof.
\end{proof}

The first property we need of~$P$ is that it is a projective sub-measurement.
This is shown in the following lemma.

\begin{lemma}[Projectivity of~$P$]\label{lem:P-projectivity}
$P = \{P_a\}$ forms a projective sub-measurement.
\end{lemma}
\begin{proof}
Let $a, b$ be (possibly distinct) outcomes.
Then
\begin{align*}
P_a \cdot P_b
&= (\widehat{X}^\dagger \cdot T_a \cdot \widehat{X})
	\cdot (\widehat{X}^\dagger \cdot T_b \cdot \widehat{X}) \tag{by \Cref{lem:pa-restated}}\\
&= (\widehat{X}^\dagger \cdot T_a \cdot I_{m \times m} \cdot T_b \cdot \widehat{X}) \tag{by \Cref{lem:X-hat-squared}}\\
&= (\widehat{X}^\dagger \cdot T_a  \cdot \widehat{X}) \cdot \bone[a = b] \tag{because~$T$ is a projective measurement}\\
&= P_a  \cdot \bone[a=b]. \tag{by \Cref{lem:pa-restated}}
\end{align*}
This completes the proof.
\end{proof}

The second property we need of~$P$ is that it is close to~$A$.
We will first show that it is close to~$Q$.

\begin{lemma}[$P$ is close to~$Q$]\label{lem:P-Q-approx}
\begin{equation*}
Q_a \otimes I \approx_{30\zeta^{1/4}} P_a \otimes I.
\end{equation*}
\end{lemma}
\begin{proof}
Our goal is to upper-bound the quantity
\begin{align}
&\sum_a \bra{\psi} (Q_a - P_a)^2 \otimes I \ket{\psi}\nonumber\\
=~& \sum_a \bra{\psi} (Q_a)^2 \otimes I \ket{\psi} + \sum_a \bra{\psi} (P_a)^2 \otimes I \ket{\psi} - \sum_a \bra{\psi} Q_a P_a \otimes I \ket{\psi}
	- \sum_a \bra{\psi} P_a Q_a \otimes I \ket{\psi}\nonumber\\
=~& \bra{\psi} Q \otimes I \ket{\psi} +  \bra{\psi} P\otimes I \ket{\psi} - \sum_a \bra{\psi} Q_a P_a \otimes I \ket{\psi}
	- \sum_a \bra{\psi} P_a Q_a \otimes I \ket{\psi},\label{eq:P-Q-thing-to-bound}
\end{align}
where the last step uses the projectivity of~$Q$ and~$P$.
We bound the four terms in \Cref{eq:P-Q-thing-to-bound} separately.
First, by \Cref{lem:projective-low-rank-sum},
\begin{equation*}
 \bra{\psi} Q \otimes I \ket{\psi}
 \leq (1+2\sqrt{\zeta}) \cdot \bra{\psi} I \otimes I \ket{\psi} \leq 1+2\sqrt{\zeta}.
\end{equation*}
Second, because~$P$ is a sub-measurement,
\begin{equation*}
\bra{\psi} P \ot I\ket{\psi} \leq 1.
\end{equation*}

The third term and fourth terms are significantly more complicated to bound.
As they are complex conjugates of each other, we can write their sum as
\begin{equation}\label{eq:complex-conjugates}
\sum_a \bra{\psi} Q_a P_a \otimes I \ket{\psi}
	+ \sum_a \bra{\psi} P_a Q_a \otimes I \ket{\psi}
= 2 \cdot \mathfrak{R}\Big(\sum_a \bra{\psi} Q_a P_a \otimes I \ket{\psi}\Big).
\end{equation}
We now focus on the expression on the right-hand side of \Cref{eq:complex-conjugates}.
To begin, we use \Cref{lem:qa-restated} to rewrite it as
\begin{equation*}
\sum_a \bra{\psi} Q_a P_a \otimes I \ket{\psi}
= \sum_a \bra{\psi} ((X_a^\dagger \cdot X) \cdot P_a)  \otimes I \ket{\psi}.
\end{equation*}
The main step will be to show that we can exchange the second $X$ for an $\widehat{X}$, i.e.
\begin{equation}\label{eq:add-a-hat}
\sum_a \bra{\psi} (X_a^\dagger \cdot X \cdot P_a) \otimes I \ket{\psi}
\approx_{2\zeta^{1/4}} \sum_a \bra{\psi} (X_a^\dagger \cdot \widehat{X} \cdot P_a)  \otimes I \ket{\psi}.
\end{equation}
To show this, we bound the magnitude of the difference.
\begin{multline*}
\Big|\sum_a \bra{\psi} ((X_a^\dagger \cdot (X - \widehat{X})) \ot I) \cdot(P_a \ot I) \ket{\psi}\Big|\\
\leq \sqrt{\sum_a \bra{\psi}(X_a^\dagger \cdot (X - \widehat{X})  \cdot (X - \widehat{X})^\dagger \cdot X_a) \ot I \ket{\psi}}\cdot \sqrt{\sum_a \bra{\psi} (P_a)^2 \ot I \ket{\psi}}.
\end{multline*}
The expression inside the first square root is
\begin{align*}
& \sum_a \bra{\psi}(X_a^\dagger \cdot (X - \widehat{X})  \cdot (X - \widehat{X})^\dagger \cdot X_a) \ot I \ket{\psi}\\
\leq~&  \sum_a \bra{\psi} (X_a^\dagger \cdot (X \cdot X^\dagger - I_{m \times m})^2 \cdot X_a) \ot I \ket{\psi}\tag{by \Cref{lem:squared-difference}}\\
=~&  \sum_a \bra{\psi} (Q_a \cdot Q \cdot Q_a - Q_a)  \ot I \ket{\psi}\tag{by \Cref{lem:X-expression-to-Q-expression}}\\
\leq~& 4\sqrt{\zeta} \cdot \bra{\psi} I \ot I \ket{\psi} \tag{by~\Cref{lem:q-almost-projective}}\\
=~& 4\sqrt{\zeta}.
\end{align*}
The expression inside the second square root is at most~$1$ because~$P$ is a sub-measurement.
Next, we claim that the \Cref{eq:add-a-hat} is in fact a real-valued expression.
To see this,
\begin{align*}
\eqref{eq:add-a-hat}
& = \sum_a \bra{\psi} (X_a^\dagger \cdot \widehat{X} \cdot P_a)  \otimes I \ket{\psi}\\
&= \sum_a \bra{\psi} (X^\dagger \cdot T_a \cdot \widehat{X} \cdot \widehat{X}^\dagger \cdot T_a \cdot \widehat{X}) \ot I \ket{\psi}
			\tag{by \Cref{lem:xa-t,lem:pa-restated}}\\
&= \sum_a \bra{\psi} (X^\dagger \cdot T_a \cdot I_{m \times m} \cdot T_a \cdot \widehat{X}) \ot I \ket{\psi} \tag{by \Cref{lem:X-hat-squared}}\\
&= \sum_a \bra{\psi} (X^\dagger \cdot T_a  \cdot \widehat{X}) \ot I \ket{\psi}\\
&=  \bra{\psi} (X^\dagger  \cdot \widehat{X}) \ot I \ket{\psi} \tag{because~$T$ is a measurement}\\
&=  \bra{\psi} \sqrt{Q} \ot I \ket{\psi} \tag{by \Cref{lem:X-times-X-hat}},
\end{align*}
which is real-valued because~$Q$ is positive semidefinite.
As a result, we have
\begin{align*}
\mathfrak{R}\Big(\sum_a \bra{\psi} Q_a P_a \otimes I \ket{\psi}\Big)
& \geq \mathfrak{R}\Big(\sum_a \bra{\psi} (X_a^\dagger \cdot \widehat{X} \cdot P_a)  \otimes I \ket{\psi}\Big) - 2\zeta^{1/4} \tag{by \Cref{eq:add-a-hat}}\\
& = \bra{\psi} \sqrt{Q} \ot I \ket{\psi}- 2\zeta^{1/4}\\
& \geq (1 - 12 \zeta^{1/4}) - 2\zeta^{1/4}. \tag{by \Cref{lem:sqrt-Q-completeness}}
\end{align*}
In total, \Cref{eq:complex-conjugates} shows that
\begin{equation*}
\sum_a \bra{\psi} Q_a P_a \otimes I \ket{\psi}
	+ \sum_a \bra{\psi} P_a Q_a \otimes I \ket{\psi}
\geq 2 \cdot(1 - 14\zeta^{1/4}).
\end{equation*}
Putting everything together, we conclude that
\begin{equation*}
\eqref{eq:P-Q-thing-to-bound}
\leq (1 + 2\sqrt{\zeta}) + 1 - 2\cdot (1-14\zeta^{1/4}) 
= 2\sqrt{\zeta} + 28 \zeta^{1/4}
\leq 30 \zeta^{1/4}.
\end{equation*}
This completes the proof.
\end{proof}

Finally, we have that
\begin{align*}
A_a \ot I
&\approx_{12 \sqrt{\zeta}} Q_a \ot I \tag{by \Cref{lem:projective-low-rank-sum}}\\
&\approx_{30 \zeta^{1/4}} P_a \ot I. \tag{by \Cref{lem:P-Q-approx}}
\end{align*}
Hence, \Cref{prop:triangle-inequality-for-approx_delta} implies that
\begin{equation*}
A_a \ot I \approx_{84 \zeta^{1/4}} P_a \ot I.
\end{equation*}
This completes the proof.
\end{proof}